\section{Đặt vấn đề}
\label{sec:dat-van-de}

\subsection{Sự phát triển của AI tại Edge}

Các hệ thống thị giác máy tính trước đây thường truyền dữ liệu từ camera về
máy chủ hoặc trung tâm dữ liệu để xử lý. Cách tiếp cận này thuận tiện cho việc
tập trung tài nguyên nhưng tạo ra độ trễ truyền thông, tiêu thụ băng thông và làm
tăng rủi ro lộ dữ liệu. Edge AI đưa mô hình suy luận đến gần nguồn sinh dữ liệu,
cho phép thiết bị phản hồi ngay cả khi kết nối mạng không ổn định và chỉ gửi về
đám mây các kết quả đã được chọn lọc.

Sự phổ biến của GPU nhúng, NPU và các bộ thực thi tối ưu đã mở rộng phạm vi ứng
dụng Edge AI trong giám sát an toàn, giao thông, nông nghiệp, bán lẻ và tự động
hóa công nghiệp. Trong các ứng dụng này, mô hình phát hiện đối tượng thuộc họ
YOLO phù hợp nhờ khả năng xử lý một giai đoạn và cân bằng tương đối tốt giữa độ
chính xác với tốc độ \cite{redmon2016,bochkovskiy2020}. Tuy nhiên, một mô hình
chạy được trên máy trạm chưa đồng nghĩa với việc có thể vận hành ổn định trên
thiết bị biên.

\subsection{Hạn chế của triển khai AI trên Edge Device}

Thiết bị biên bị giới hạn đồng thời bởi nhiều loại tài nguyên. Các giới hạn này
có quan hệ với nhau và phải được xem xét trên toàn bộ chuỗi suy luận thay vì chỉ
đo thời gian thực thi của riêng mạng nơ-ron.

\begin{itemize}
  \item \textbf{Tài nguyên GPU:} số nhân xử lý và thông lượng tính toán thấp hơn
        máy chủ, trong khi GPU còn có thể được chia sẻ cho giải mã video, hiển
        thị và các tác vụ khác.
  \item \textbf{Bộ nhớ:} dung lượng RAM và bộ nhớ GPU giới hạn kích thước mô
        hình, batch size, số luồng camera và lượng bộ đệm có thể duy trì.
  \item \textbf{Điện năng:} thiết bị thường hoạt động theo power mode cố định;
        tăng xung nhịp có thể cải thiện FPS nhưng làm tăng công suất tiêu thụ.
  \item \textbf{Nhiệt độ:} tải suy luận kéo dài có thể làm thiết bị giảm xung để
        bảo vệ phần cứng, khiến hiệu năng thực tế thấp hơn phép đo ngắn hạn.
  \item \textbf{Yêu cầu thời gian thực:} hệ thống phải kiểm soát cả độ trễ trung
        bình, độ trễ phần vị cao và độ trễ đầu cuối từ camera đến kết quả.
\end{itemize}

Vì vậy, bài toán triển khai cần tối ưu đồng thời độ chính xác, tốc độ, kích thước
mô hình, bộ nhớ và điện năng. Mọi so sánh phải được thực hiện trên cùng dữ liệu,
cùng kích thước đầu vào, cùng thiết bị và cùng quy trình đo.

\subsection{Bài toán quản lý và triển khai AI Model trên nhiều Edge Device}

Khi số lượng thiết bị tăng, việc sao chép thủ công tệp mô hình và khởi động tiến
trình qua từng phiên SSH không còn phù hợp. Người vận hành cần biết thiết bị nào
đang trực tuyến, phiên bản mô hình đang chạy, cấu hình mong muốn, trạng thái
thực tế và lịch sử thay đổi. Một lỗi triển khai phải được phát hiện, khoanh vùng và
khôi phục mà không làm gián đoạn toàn bộ đội thiết bị.

Bài toán vì thế bao gồm cả mặt phẳng dữ liệu tại Edge và mặt phẳng điều khiển
trên Cloud. Cloud Control Plane quản lý thiết bị, runtime, deployment, lệnh điều
khiển và dữ liệu giám sát. Edge Agent tiếp nhận desired state, tải model, điều khiển
AI Runtime và báo cáo trạng thái, telemetry. Mỗi thao tác phải có định danh, trạng
thái và lịch sử để có thể truy nguyên khi xảy ra sự cố.

\subsection{Sự cần thiết của tối ưu mô hình AI}

Mua phần cứng mạnh hơn chỉ giải quyết một phần vấn đề và làm tăng chi phí khi
quy mô đội thiết bị mở rộng. Tối ưu mô hình giúp khai thác hiệu quả tài nguyên
hiện có, tăng số luồng camera có thể phục vụ và giảm thời gian phản hồi. Trong đồ
án này, pruning được dùng để loại bỏ phần tham số hoặc cấu trúc ít quan trọng;
Quantization-Aware Training (QAT) giúp mô hình thích nghi với sai số lượng tử;
TensorRT chuyển đổi đồ thị sang engine tối ưu cho GPU NVIDIA.

Tối ưu chỉ có ý nghĩa khi mức giảm chi phí suy luận không làm chất lượng nhận
diện suy giảm vượt ngưỡng chấp nhận. Do đó, đồ án không tìm mô hình nhỏ nhất
hoặc nhanh nhất một cách độc lập mà tìm cấu hình có trade-off phù hợp giữa
accuracy, latency, FPS, model size, GPU memory và power trên Jetson Nano và
Jetson Orin Nano.